\documentclass[aspectratio=169,11pt]{beamer}

\usetheme{Madrid}
\usecolortheme{default}
\usefonttheme{professionalfonts}
\setbeamertemplate{navigation symbols}{}
\setbeamertemplate{footline}[frame number]

\usepackage{amsmath, amssymb, booktabs, graphicx}

\title{Institutional Persistence, Path Dependence, and Historical Labor-Market Inequality}
\subtitle{Labor Markets and Political/Cultural Institutions -- Week 7}
\author{Teaching deck draft}
\date{}

\begin{document}

\begin{frame}
  \titlepage
\end{frame}

\begin{frame}{Central question}
\begin{itemize}
    \item Why do past labor institutions shape modern labor outcomes after the original rules disappear?
    \item Which mechanisms connect historical exposure to wages, mobility, occupations, skills, and bargaining power?
    \item When is a long-run pattern persistence, and when is it just serial correlation in fundamentals?
\end{itemize}
\end{frame}

\begin{frame}{Course reorientation}
\begin{itemize}
    \item Week 6: current hierarchy as a system of access to jobs, networks, credentials, and protection.
    \item Week 7: historical institutions as origins of durable labor-market inequality.
    \item Week 8: reform and implementation when inherited structures shape adjustment.
    \item The course object remains labor allocation, not generic historical political economy.
\end{itemize}
\end{frame}

\begin{frame}{Why history belongs in labor economics}
\begin{itemize}
    \item Labor institutions set outside options, mobility costs, bargaining power, and contract terms.
    \item Historical rules shape schools, roads, land, local public goods, and market access.
    \item Networks and political rights transmit job information and worker protection over time.
    \item The modern outcome must be a labor margin: wages, jobs, occupations, mobility, skills, or power.
\end{itemize}
\end{frame}

\begin{frame}{What counts as persistence?}
\small
\begin{tabular}{p{0.29\linewidth} p{0.31\linewidth} p{0.29\linewidth}}
\toprule
Concept & Mechanism & Labor diagnostic \\
\midrule
Persistent fundamentals & geography, resources, transport, disease, urban potential & Does history proxy a stable baseline? \\
Path dependence & increasing returns, sunk costs, networks, lock-in & Did a shock move the economy onto a durable trajectory? \\
Institutional persistence & old rules change current labor channels & Which labor margin is still moving? \\
\bottomrule
\end{tabular}
\end{frame}

\begin{frame}{A persistence decomposition}
\[
Y_{it}^{labor} = \alpha + \beta H_i + \delta M_{it} + \gamma X_i + \lambda_t + \varepsilon_{it}
\]
\begin{itemize}
    \item $H_i$: historical institutional exposure.
    \item $M_{it}$: measured labor mechanism: schooling, mobility, infrastructure, employer power, networks.
    \item $X_i$: persistent fundamentals.
    \item Goal: interpret $\beta$ through a labor mechanism, not only a long-run coefficient.
\end{itemize}
\end{frame}

\begin{frame}{Core warning}
\begin{itemize}
    \item A historical coefficient is not yet a persistence mechanism.
    \item A strong paper names the old institution, the counterfactual, and the first moving labor margin.
    \item It also asks which stable fundamentals could produce the same pattern.
    \item The question to keep asking: what does the old institution still change about work today?
\end{itemize}
\end{frame}

\begin{frame}{Coercion, labor power, and contract enforcement}
\begin{itemize}
    \item Coercive systems restrict mobility and weaken outside options.
    \item Dell: Peru's mining mita boundary as historical exposure to forced labor.
    \item Key channels: public goods, roads, market access, land, local elite power.
    \item Naidu--Yuchtman: coercive contract enforcement changes quitting costs, wages, and adjustment.
\end{itemize}
\end{frame}

\begin{frame}{Coercion to modern labor outcomes}
\small
\begin{tabular}{p{0.30\linewidth} p{0.30\linewidth} p{0.30\linewidth}}
\toprule
Historical object & First labor margin & Modern outcomes \\
\midrule
Labor draft & mobility and outside options & wages, occupations, self-employment \\
Coercive contract law & quitting costs and discipline & wages, separations, turnover \\
Serfdom or slavery & land access and bargaining & migration, schooling, employer power \\
\bottomrule
\end{tabular}
\end{frame}

\begin{frame}{Migration regimes and mobility}
\begin{itemize}
    \item Migration restrictions shape settlement, job networks, and local information.
    \item Forced migration can change local public finance, housing, occupations, and politics.
    \item The labor question is not only who moved, but which labor markets became accessible.
    \item Diagnostic margins: commuting, migration, entrepreneurship, occupation sorting, job entry.
\end{itemize}
\end{frame}

\begin{frame}{Schooling, public goods, and occupational structure}
\begin{itemize}
    \item Land and local public finance affect schools, roads, market access, and labor demand.
    \item Markevich--Zhuravskaya: serfdom abolition changes mobility, competition, and local structure.
    \item Jones--Schmick: Reconstruction education and long-run Black-White labor inequality.
    \item The frontier ties historical institutions to measured schooling and occupational standing.
\end{itemize}
\end{frame}

\begin{frame}{Networks, information, and political rights}
\begin{itemize}
    \item Networks transmit job information, insurance, and exit costs.
    \item Political rights affect school finance, public employment, enforcement, and protection.
    \item Community institutions can support workers while locking in low-return occupations.
    \item Persistence can survive formal abolition through these meso-level channels.
\end{itemize}
\end{frame}

\begin{frame}{Wars and abrupt shocks}
\begin{itemize}
    \item War shocks can shift labor demand, migration, skill acquisition, discrimination, and politics.
    \item Aizer--Boone--Lleras-Muney--Vogel: WWII and racial disparities in labor outcomes.
    \item Ferrara: wars and minority labor-market outcomes in the twentieth century.
    \item Main risk: war bundles many treatments, so the labor margin must be narrowed.
\end{itemize}
\end{frame}

\begin{frame}{Methods box: historical designs}
\small
\begin{tabular}{p{0.31\linewidth} p{0.30\linewidth} p{0.29\linewidth}}
\toprule
Design & Identifies best & Main caution \\
\midrule
Historical border or discontinuity & local exposure to an old institution & geography and boundary selection \\
War or abrupt political shock & demand, migration, or access shifts & bundled treatment \\
Regime abolition or reform & removal or expansion of institutions & simultaneous reforms \\
Linked microdata & lifecycle and intergenerational effects & linkage and survival bias \\
Lineage or ancestry proxy & family and community transmission & selection and mixed mechanisms \\
Historical GIS & spatial access and public goods & spatial correlation and map error \\
\bottomrule
\end{tabular}
\end{frame}

\begin{frame}{Data box}
\begin{itemize}
    \item Historical census and linked census data: occupation, literacy, migration, household structure.
    \item Archives, courts, contracts, and payrolls: enforcement, breach, quitting, wages, personnel.
    \item School, parish, and public-goods records: teachers, finance, roads, hospitals, poor relief.
    \item Cadastral and colonial maps, GIS layers, war records, military files, newspapers, legal texts.
\end{itemize}
\end{frame}

\begin{frame}{Representative paper map}
\small
\begin{tabular}{p{0.27\linewidth} p{0.33\linewidth} p{0.30\linewidth}}
\toprule
Theme & Anchor & Labor mechanism \\
\midrule
Forced labor & Dell & public goods, roads, market access \\
Contract coercion & Naidu--Yuchtman & quitting costs, wages, mobility \\
Abolition & Markevich--Zhuravskaya & competition, mobility, organization \\
Education and race & Jones--Schmick & schooling, occupation, inequality \\
War reallocation & Aizer et al.; Ferrara & demand, discrimination, openings \\
Forced migration & Chevalier et al. & settlement, policy, local labor markets \\
\bottomrule
\end{tabular}
\end{frame}

\begin{frame}{Research frontier}
\begin{itemize}
    \item Measure modern employer power as a legacy of coercive labor systems.
    \item Link historical schooling and public goods to occupation and earnings trajectories.
    \item Use digitized archives, OCR, linked census files, and historical GIS to test mechanisms.
    \item Study when forced migration changes labor protection, local demand, or sorting.
    \item Convert broad legacy claims into concrete labor-allocation mechanisms.
\end{itemize}
\end{frame}

\begin{frame}{Week 7 lab}
\begin{itemize}
    \item Reproduce: synthetic Dell-style boundary comparison.
    \item Diagnose: classify persistence claims as fundamentals, path dependence, or labor-channel persistence.
    \item Transfer: map designs to Naidu--Yuchtman, Markevich--Zhuravskaya, Jones--Schmick, wars, and forced migration.
    \item Output: reproduction table, mechanism diagnostic, transfer classification.
\end{itemize}
\end{frame}

\begin{frame}{Bridge to Week 8}
\begin{itemize}
    \item Week 7 explains why inherited institutions can survive in labor markets.
    \item Week 8 asks how reforms work when those inherited structures differ.
    \item Reform effects depend on employer power, mobility costs, skills, public goods, networks, rights, and capacity.
    \item The next question: when does new law change actual labor-market behavior?
\end{itemize}
\end{frame}

\begin{frame}{Takeaways}
\begin{itemize}
    \item Persistence is a mechanism claim, not only a time-series fact.
    \item Historical labor institutions persist through outside options, skills, mobility, public goods, networks, and rights.
    \item Identification designs are useful only when paired with a labor mechanism and careful data provenance.
    \item The frontier is historical, but the object remains work, wages, allocation, power, and inequality.
\end{itemize}
\end{frame}

\end{document}
